\chapter{Introduction to Spatiotemporal Weighted Value Distribution Maps}\label{chap:christian-introduction}

This chapter provides an overview of the desired result of this part of the thesis, including its requirements and terminology. It also examines existing solutions.

\section{Overview of input data and desired result}

For the purpose of visualizing the sensor data, one data type is relevant: a \emph{measurement}, which represents a single data point with the following information:
\begin{itemize}
    \item Time of measurement and upload to the servers
    \item Geographical position in latitude and longitude
    \item The user responsible for this measurement
    \item Humidity
    \item Temperature
    \item Fine dust concentration for particles of 1$\mu m$, 2.5$\mu m$, and 10$\mu m$ in size
    \item Air pressure
\end{itemize}

The goal is to create a map similar to a heat map, displaying an overlay
derived from data points on a map.


\autoref{fig:Map_01} shows a heat map of the frequency and location of
Instagram photos tagged \#aokhalifax~\cite{heatmap:01}. Warmer colors represent
a higher \emph{density} of data points, in this case Instagram posts.

Instead of density, our map should display the \emph{associated value} of the data points and interpolate
regions without data points.

% \begin{figure}[H]
%     \centering
%     \includegraphics[scale=0.38]{images/valuedistributionmap.jpg}
%     \caption{A manually drawn example of the desired result, original map from OpenStreetMap~\cite{openstreetmap:01}}\label{fig:Map_02}
% \end{figure}

\autoref{fig:Map_02} shows a hand-drawn depiction of the desired result. It
shows an arbitrary quantity over St.\ Pölten, ranging from white to brown. The density of data points can only be discerned through the opacity of the overlay.

The final requirements of the map are:
\begin{itemize}
    \item The opacity of the overlay corresponds to the \emph{density} of data points
    \item The color of the overlay is an interpolation of the \emph{associated value}
          from nearby data points.
    \item It should be reasonably performant and preferably run on mobile and low-end devices.
\end{itemize}

\section{Terminology and Naming}

An important task was finding terminology to describe a map with the above
requirements. This is necessary to conduct research on the topic and establish
a clear title for this thesis. According to ChatGPT~\cite{chatgpt:01}, the map
that is similar to a heat map but displays value instead of density is called a
\emph{spatiotemporal weighted value distribution map}. Upon further
investigation, this was revealed to only be a description of the goals, made up
of the following parts:
\begin{itemize}
    \item Spatiotemporal: location and time data are being processed
    \item Weighted: data points can be weighted, adjusting for age and possibly clumping
          of similar data for performance reasons
    \item Value distribution: the map shows the distribution of a certain value, and not
          density or counts
    \item Map: the resulting visualization shows that value over space
\end{itemize}

ChatGPT also provided some formally recognized terms:
\begin{itemize}
    \item Spatiotemporal Kernel Regression: from a statistical standpoint
    \item Spatiotemporal Interpolation: from a geographic standpoint
\end{itemize}
All of these fall under the umbrella term \emph{spatiotemporal interpolation methods}, or \emph{spatiotemporal smoothing methods}~\cite{mdpi:02}. This leads to the first part of the title of this thesis: \emph{Theory of Weighted Value Distribution Maps Including the Comparison of Spatiotemporal Smoothing Methods}, which is covered in \autoref{ch:theory}.

In this thesis, it will simply be referred to as \emph{the map}.

\section{Existing Solutions}

\subsection{Surfer}
Surfer is a program developed by Golden Software~\cite{surfer:01} for
the visualization of geographic data. It offers ``numerous interpolation methods to
grid regularly or irregularly spaced data onto a grid or raster''\cite{surfer:02}, such as kriging, IDP\footnote{Inverse Distance to a Power} in
two or three dimensions. The different interpolation methods will be discussed
later (see \autoref{sec:smoothingmethods}). In addition, Surfer supports a great
number of tools for computational analysis, placing data on a map, and the
analysis of three-dimensional data.

Surfer fulfills the requirements of the map; however, it is paid software, and
the goal is to have a dynamic map on a website, rendered in a browser, so it
cannot be used. However, this thesis will use some of the interpolation methods
mentioned in the Surfer specifications.

\subsection{The gstat R package}
\verb|gstat|~\cite{rgstat:paper1, rgstat:paper2} is an R package used for interpolation of geospatial data. It has support for both spatial and spatiotemporal modeling, making it a good choice for producing value distribution maps such as the one in this thesis.

\verb|gstat| is provided with data that include location and time attributes. The package then estimates how values correlate across distance and time using a variogram model. Once that relationship is defined, \verb|gstat| performs Kriging (see~\autoref{subsec:kriging}) to predict values at unmeasured locations and times. The output is a continuous surface or series of surfaces that can be visualized as maps or animations. It works alongside data structures from multiple different packages, such as \verb|sp|\footnote{Spatial Data} and \verb|sf|\footnote{Simple features} for spatial vector data, or \verb|raster| and \verb|stars| for spatial raster data, allowing the final distribution to be used with standard R mapping workflows.\ \cite{r:spatialdata}

\section{Comparison with the desired Result}

All of the solutions mentioned above fail to meet at least one of the
requirements. The main problem is the absence of automation (in software such as
Surfer), which is necessary for the project. Additionally, the ability to
display results on the client-side using JavaScript is not an option with
existing solutions.

Although using R for data processing is a potential option, it will not be
considered in this thesis due to the specific goals of the project. The focus
of this thesis is the development of a custom system that can meet the unique
needs of the project.